\documentclass[11pt,a4paper]{article}

%========================
% PACKAGES
%========================
\usepackage[utf8]{inputenc}
\usepackage[T1]{fontenc}
\usepackage[french]{babel}
\usepackage{lmodern}
\usepackage{amsmath,amssymb}
\usepackage{geometry}
\usepackage{enumitem}
\usepackage{xcolor}
\usepackage{tcolorbox}
\usepackage{fancyhdr}
\usepackage{array}
\usepackage{multicol}

\geometry{margin=2cm}

%========================
% COULEURS
%========================
\definecolor{bleu}{RGB}{25,70,130}
\definecolor{grisclair}{RGB}{245,247,250}
\definecolor{rouge}{RGB}{160,30,30}
\definecolor{vert}{RGB}{20,120,70}

%========================
% EN-TÊTE
%========================
\pagestyle{fancy}
\fancyhf{}
\lhead{\textbf{Mathématiques — 6e}}
\rhead{\textbf{Devoir maison}}
\cfoot{\thepage}

%========================
% BOÎTES
%========================
\newtcolorbox{consigne}{
  colback=grisclair,
  colframe=bleu,
  boxrule=0.8pt,
  arc=2mm,
  left=2mm,
  right=2mm,
  top=1mm,
  bottom=1mm
}

\newtcolorbox{attention}{
  colback=red!4!white,
  colframe=rouge,
  boxrule=0.8pt,
  arc=2mm,
  left=2mm,
  right=2mm,
  top=1mm,
  bottom=1mm
}

%========================
% DOCUMENT
%========================
\begin{document}

\begin{center}
    {\Large \textbf{\textcolor{bleu}{Devoir maison de mathématiques}}}\\[2mm]
    {\large Niveau 6e — Raisonnement, calcul et logique}\\[2mm]
    \textbf{Nom :} \dotfill \qquad
    \textbf{Prénom :} \dotfill \\[2mm]
    \textbf{Date :} \dotfill
\end{center}

\vspace{3mm}

\begin{consigne}
\textbf{Objectif du devoir :} ce devoir ne cherche pas seulement à vérifier que tu sais calculer.
Il cherche surtout à voir si tu sais \textbf{raisonner}, \textbf{expliquer tes choix}, \textbf{organiser tes calculs}
et \textbf{rédiger une réponse claire}.
\end{consigne}

\vspace{2mm}

\begin{attention}
\textbf{Consignes importantes :}
\begin{itemize}
    \item Tous les calculs doivent être posés ou expliqués.
    \item Une réponse sans justification peut rapporter peu de points, même si le résultat est juste.
    \item Les phrases de conclusion sont attendues dans les problèmes.
    \item Il est autorisé de chercher, de se tromper, puis de corriger proprement.
\end{itemize}
\end{attention}

\vspace{4mm}

\begin{center}
\renewcommand{\arraystretch}{1.3}
\begin{tabular}{|c|c|}
\hline
\textbf{Exercice} & \textbf{Barème indicatif} \\
\hline
Exercice 1 & 4 points \\
Exercice 2 & 5 points \\
Exercice 3 & 5 points \\
Exercice 4 & 6 points \\
Exercice 5 & 6 points \\
Exercice 6 & 4 points \\
\hline
\textbf{Total} & \textbf{30 points} \\
\hline
\end{tabular}
\end{center}

\newpage

%========================
\section*{Exercice 1 — Calculs réfléchis et explications \hfill /4}

Pour chaque question, donne le résultat puis explique ta méthode.

\begin{enumerate}[label=\textbf{\arabic*)}]
    \item Calcule mentalement, en expliquant l'astuce utilisée :
    \[
    199 + 48 + 101 + 52.
    \]

    \item Calcule sans poser la multiplication :
    \[
    25 \times 36.
    \]
    Tu dois proposer une méthode utilisant une décomposition de 36.

    \item Complète les égalités suivantes :
    \[
    7 \times \Box = 84
    \qquad
    \Box \times 9 = 108
    \qquad
    144 \div \Box = 12.
    \]

    \item Léa affirme :
    \[
    38 \times 12 = 38 \times 10 + 38 \times 2.
    \]
    Explique pourquoi son raisonnement est correct, puis termine le calcul.
\end{enumerate}

\vspace{4mm}

%========================
\section*{Exercice 2 — Problème de logique : le nombre mystère \hfill /5}

Je pense à un nombre entier.

Voici les indices :

\begin{itemize}
    \item Il est plus grand que 40 et plus petit que 90.
    \item Il est divisible par 3.
    \item Il est pair.
    \item La somme de ses chiffres est égale à 12.
\end{itemize}

\begin{enumerate}[label=\textbf{\arabic*)}]
    \item Écris tous les nombres pairs compris entre 40 et 90 dont la somme des chiffres est 12.
    \item Parmi ces nombres, lesquels sont divisibles par 3 ?
    \item Quel est ou quels sont les nombres possibles ?
    \item On ajoute l'indice suivant :
    \begin{center}
    \og Le nombre est plus proche de 80 que de 60. \fg
    \end{center}
    Quel est alors le nombre mystère ?
    \item Rédige une phrase expliquant ton raisonnement.
\end{enumerate}

\vspace{4mm}

%========================
\section*{Exercice 3 — Organiser un raisonnement \hfill /5}

Une école organise une sortie au musée.

Il y a :
\begin{itemize}
    \item 4 classes de 6e ;
    \item 26 élèves dans chaque classe ;
    \item 7 accompagnateurs ;
    \item des bus de 55 places chacun.
\end{itemize}

\begin{enumerate}[label=\textbf{\arabic*)}]
    \item Combien y a-t-il d'élèves au total ?
    \item Combien de personnes participent à la sortie en comptant les accompagnateurs ?
    \item Combien de bus faut-il réserver au minimum ?
    \item Combien restera-t-il de places libres dans les bus ?
    \item Un élève écrit :
    \[
    4 \times 26 + 7 \div 55.
    \]
    Explique pourquoi ce calcul ne permet pas de répondre directement au problème.
\end{enumerate}

\vspace{4mm}

%========================
\section*{Exercice 4 — Géométrie et précision \hfill /6}

On considère un rectangle $ABCD$ tel que :
\[
AB = 8 \text{ cm}
\qquad \text{et} \qquad
BC = 5 \text{ cm}.
\]

\begin{enumerate}[label=\textbf{\arabic*)}]
    \item Construis le rectangle $ABCD$ avec soin sur ta copie.
    \item Calcule le périmètre du rectangle.
    \item Calcule l'aire du rectangle.
    \item On place un point $E$ sur le segment $[AB]$ tel que :
    \[
    AE = 3 \text{ cm}.
    \]
    Construis le point $E$.
    \item Calcule la longueur $EB$.
    \item Trace le segment $[EC]$. Explique pourquoi la figure obtenue contient deux triangles.
    \item Question de réflexion : les deux triangles formés ont-ils forcément la même aire ? Justifie ta réponse par un calcul ou une explication claire.
\end{enumerate}

\vspace{4mm}

%========================
\section*{Exercice 5 — Un problème à plusieurs étapes \hfill /6}

Un marchand vend des cahiers et des stylos.

\begin{itemize}
    \item Un cahier coûte 3 euros.
    \item Un stylo coûte 2 euros.
    \item Malik achète 5 cahiers et 4 stylos.
    \item Sarah achète 3 cahiers et 8 stylos.
\end{itemize}

\begin{enumerate}[label=\textbf{\arabic*)}]
    \item Calcule le prix payé par Malik.
    \item Calcule le prix payé par Sarah.
    \item Qui a payé le plus cher ? De combien ?
    \item Le marchand propose ensuite une réduction de 5 euros si le total dépasse 25 euros.
    
    Malik bénéficie-t-il de la réduction ? Justifie.
    
    Sarah bénéficie-t-elle de la réduction ? Justifie.

    \item Après réduction, combien paie chacun ?
    
    \item Question de raisonnement :
    
    Thomas veut acheter seulement des cahiers et des stylos. Il veut dépenser exactement 20 euros.
    
    Trouve au moins trois achats différents possibles.
    
    Présente tes réponses dans un tableau.
\end{enumerate}

\vspace{4mm}

%========================
\section*{Exercice 6 — Vrai ou faux ? Justifie \hfill /4}

Pour chaque affirmation, écris \textbf{Vrai} ou \textbf{Faux}, puis justifie ta réponse.

\begin{enumerate}[label=\textbf{\arabic*)}]
    \item Si un nombre est pair, alors il est divisible par 4.
    
    \item Si un nombre est divisible par 10, alors il est pair.
    
    \item Le périmètre d'un carré de côté 6 cm est égal à son aire.
    
    \item Deux rectangles qui ont le même périmètre ont forcément la même aire.
    
    \item Quand on multiplie un nombre par 3 puis par 2, c'est pareil que le multiplier par 6.
    
    \item Le nombre 1 est un multiple de tous les nombres entiers.
    
    \item Si on ajoute 100 à un nombre, son chiffre des unités ne change pas.
    
    \item Un triangle peut avoir trois côtés de même longueur.
\end{enumerate}

\vspace{5mm}

\begin{center}
\Large \textbf{\textcolor{bleu}{Fin du devoir}}
\end{center}

\vspace{3mm}

\begin{consigne}
\textbf{Conseil final :} relis ton devoir. Vérifie que chaque réponse importante est accompagnée d'une justification.
Un bon devoir de mathématiques n'est pas seulement une suite de résultats : c'est un raisonnement écrit clairement.
\end{consigne}

\end{document}